% Calibrated loop control for multi-hop QA with small language models (CVPR/AVSS camera-ready format).
% Build: tectonic paper/main.tex   (see paper/README.md; Overleaf fallback documented there)
%
% Paragraph markers: every English paragraph/float is preceded by a comment line
% "%% [§<section> ¶<n>]". paper/zh-TW/main.md carries the same markers as
% "<!-- EN: §<section> ¶<n> -->" so the two versions can be aligned mechanically
% (paper/check_paper.sh verifies this).
%
% The paper contains no placeholders; paper/check_paper.sh fails on any leftover marker.
% Revision R9: CVPR author kit (cvpr.sty 2026, ieeenat_fullname.bst), the layout of the reference
% paper (IEEE AVSS 2025): camera-ready with page numbers, numbered sections, captions "Table 1." above
% and "Figure 1." below, full-name references at 10 pt, and no keywords.
% Authors use the CVPR kit's separate author blocks, with an ORCID link for each author.
% R9 round 2: main text + References = exactly 8 pages, page 8 filled; then, from page 9 and outside the
% page budget, Appendix A (implementation and reproducibility details) and B (supplementary results).
\documentclass[10pt,twocolumn,letterpaper]{article}

% Camera-ready version with page numbers, as in the reference paper. cvpr.sty loads times, xcolor,
% graphicx, amsmath, amssymb, booktabs, natbib [numbers,sort&compress], caption, url, and enumitem.
\usepackage[pagenumbers]{cvpr}
% T1 encoding under both engines: Type 1 Times (NimbusRomNo9L, as in the reference paper) and
% \textquotedbl (the straight quote of \dq) under XeTeX (tectonic) as well as pdfLaTeX (Overleaf).
\usepackage[T1]{fontenc}
% cvpr.sty sets \hfuzz=30pt; restore the LaTeX default so that every Overfull \hbox is reported
% (paper/check_paper.sh fails on any).
\hfuzz=0.1pt
% The reference paper indents every paragraph, also the first one after a heading and the abstract
% (the 2026 kit starts the abstract with \noindent).
\usepackage{indentfirst}
\patchcmd{\abstract}{\noindent}{}{}{\PackageError{main}{cvpr.sty abstract patch failed}{}}

\usepackage{amsfonts}
\usepackage{algorithm}
\usepackage{algpseudocode}
% One-line conditional for Algorithm 1: \IfThen{<condition>}{<statement>}.
\algnewcommand{\IfThen}[2]{\State \algorithmicif\ #1\ \algorithmicthen\ #2}
% Blocks end by indentation (no "end if" / "end for" lines), to fit the narrow CVPR column.
\algtext*{EndIf}
\algtext*{EndFor}

\usepackage{textcomp}
% Fixed-width aligned columns of the ruled Table 3 panels (R9 round 3; >{\raggedleft\arraybackslash}p{...}).
\usepackage{array}

% The Appendix A heading must not break as "De-tails" (R9 round 2) nor be stretched across the column:
% section headings are set ragged right (one-line headings are unchanged).
\hyphenation{Details}
\makeatletter
\patchcmd{\cvprsection}{\large\bf}{\large\bf\raggedright}{}{\PackageError{main}{cvpr.sty section patch failed}{}}
\makeatother


% The CVPR kit's preamble.tex: microtype "works harder in optimizing text layout"; enabled at the end
% of writing to stay within the page limit (revision R9).
\usepackage{microtype}
% R9 round 2: paper/check_paper.sh reads the end of the References from main.aux: \label{refs-end} gives the
% page, \zsavepos{refs-end} the position of the last reference line (page 8 must be filled).
\usepackage{zref-savepos}
\AtEndEnvironment{thebibliography}{\label{refs-end}\zsavepos{refs-end}}


% ---------------------------------------------------------------------------
% Table/figure hook macros
% ---------------------------------------------------------------------------
% Visible red marker shown only when a generated table/figure file is absent.
\newcommand{\missingfile}[1]{\textcolor{red}{\textbf{[MISSING FILE: #1]}}}

% Contract with Ticket 11 (fixed by the Coordinator):
%   paper/tables/<stem>.tex  = a bare tabular only (booktabs rules allowed); no table
%                              float, no \caption, no \label.
%   paper/figures/<stem>.pdf = the figure itself.
% The macros below always build the float, caption, and label; the only difference
% between "file present" and "file absent" is the body (the file or a red marker).

% \tableorplaceholder[<font size>]{<file stem>}{<caption>} -> Table with \label{tab:<file stem>}
% The generated ruled tables carry no @{} padding; their cell padding is \tabcolsep = 2.5pt, set here
% (R9 round 3), so that every \cline and \multicolumn meets the vertical rules.
% (one column, top float). \widetableorplaceholder is the same as a two-column table* (revision
% R8: the efficiency panels). Main-text tables are cited in the main text and placed there; the
% appendix tables (R9 round 2) are cited and placed in the appendix.
\newcommand{\tableorplaceholder}[3][]{%
  \begin{table}[t]
    \centering
    \caption{#3}
    \label{tab:#2}#1\setlength{\tabcolsep}{2.5pt}%
    \IfFileExists{tables/#2.tex}{\input{tables/#2.tex}}{%
      \fbox{\parbox{0.92\columnwidth}{\centering\missingfile{\url{tables/#2.tex}}}}}
  \end{table}}
\newcommand{\widetableorplaceholder}[3][]{%
  \begin{table*}[t]
    \centering
    \caption{#3}
    \label{tab:#2}#1\setlength{\tabcolsep}{2.5pt}%
    \IfFileExists{tables/#2.tex}{\input{tables/#2.tex}}{%
      \fbox{\parbox{0.92\textwidth}{\centering\missingfile{\url{tables/#2.tex}}}}}
  \end{table*}}

% \figureorplaceholder[<width>]{<file stem>}{<caption>} -> Figure with \label{fig:<file stem>}
\newcommand{\figureorplaceholder}[3][\columnwidth]{%
  \begin{figure}[t]
    \centering
    \IfFileExists{figures/#2.pdf}{\includegraphics[width=#1]{figures/#2.pdf}}{%
      \fbox{\parbox[c][3cm][c]{0.92\columnwidth}{\centering\missingfile{\url{figures/#2.pdf}}}}}
    \caption{#3}
    \label{fig:#2}
  \end{figure}}

% \subfigureorplaceholder[<width>]{<file stem>}{<caption>} -> a panel (default one third of the
% text width) for use inside a figure* float; each panel keeps its own caption, number, and
% \label{fig:<file stem>}.
\newcommand{\subfigureorplaceholder}[3][0.32\textwidth]{%
  \begin{minipage}[t]{#1}
    \centering
    \IfFileExists{figures/#2.pdf}{\includegraphics[width=\linewidth]{figures/#2.pdf}}{%
      \fbox{\parbox[c][3cm][c]{0.92\linewidth}{\centering\missingfile{\url{figures/#2.pdf}}}}}
    \caption{#3}
    \label{fig:#2}
  \end{minipage}}

\newcommand{\cond}[1]{\textsc{#1}}
% Straight double quotes (U+0022). A literal " would become a curly quote through the TeX
% ligature mapping, so quoted text is written \dq{...}.
\newcommand{\dq}[1]{\textquotedbl{}#1\textquotedbl{}}

% The added ORCID row uses one author line (14 pt). Reuse 14 pt of the title-to-author
% gap so the author details fit without changing the text area's height or pagination.
\makeatletter
\patchcmd{\@maketitle}{\vspace*{24pt}}{\vspace*{10pt}}{}{\PackageError{main}{CVPR title spacing patch failed}{}}
\makeatother

% hyperref last, as in the CVPR kit; the reference paper's link colours (hyperref defaults:
% red internal links, green citations); no page back-references in the camera-ready version.
\usepackage[breaklinks,colorlinks]{hyperref}

\begin{document}

\title{Calibrated Loop Control for Multi-Hop QA with Small Language Models}

% Separate author blocks follow the official CVPR kit. Each author's name, affiliation,
% email, and ORCID occupy corresponding rows; ORCID links are black in print.
\author{Pin-Hung Lin\\
National Cheng Kung University\\
{\tt\small f84076029@gs.ncku.edu.tw}\\
{\small\hypersetup{urlcolor=black}\url{https://orcid.org/0009-0000-1916-7829}}
\and
Te-Lun Yang\\
National Taiwan University\\
{\tt\small d12944007@ntu.edu.tw}\\
{\small\hypersetup{urlcolor=black}\url{https://orcid.org/0000-0002-3351-1785}}}

\maketitle

% Numbers only from summary.csv and laya_eval_report.md of run main-20260930 and from
% paper/analysis/error_taxonomy.csv.
\begin{abstract}
%% [§0 ¶1]
Small language models (SLMs) often stop multi-hop question-answering loops too early,
retrieve unnecessarily, or misuse tools. Most adaptive retrieval methods rely on the
generator's signals. When should an external calibrated controller take over?
We study Laya, a compact non-autoregressive decision model with calibrated outputs. At
each hop, it predicts whether the evidence is sufficient, what to do next, and how much
evidence is missing. Our $2\times2$ factorial design uses a fixed decomposition loop
with two evidence sources and two controllers. Evidence comes from closed-book recall
or retrieval-augmented generation (RAG) with re-ranking. The generating large language
model (LLM) or Laya controls the loop. We evaluate five Qwen3.5 models (0.8B to 27B
parameters) on HotpotQA, 2WikiMultihopQA, and MuSiQue, using 200 questions per
configuration. Trained on weak labels, Laya achieves 0.9271 held-out accuracy on evidence
sufficiency. Retrieval improves exact match (EM) by up to 43.5 percentage points. Under
RAG, Laya raises EM point estimates by 1.5--9.0 points for the 0.8B and 2B models but
lowers EM by 3.5--15.5 points for the 4B--27B models. It reduces prompt tokens and hops
throughout. Its most frequent error is stopping before all gold evidence is retrieved.
The accuracy-efficiency trade-off of external control therefore depends on model size.
\end{abstract}


\section{Introduction}
\label{sec:intro}

%% [§1 ¶1]
Large language models (LLMs) can hallucinate, rely on outdated knowledge, and lack
specialised facts. Retrieval-augmented generation (RAG)~\cite{lewis2020rag} provides
the generator with passages from an external corpus. Combining dense search with a
cross-encoder re-ranker further improves question-answering (QA) accuracy across many
LLMs~\cite{yang2025enhancing}.

%% [§1 ¶2]
Multi-hop questions combine facts from multiple documents, for example to identify a
bridge entity and ask about its
properties~\cite{yang2018hotpotqa,ho2020twowiki,trivedi2022musique}. A single retrieval
often misses second-hop evidence. Agents therefore decompose questions into sub-questions
and gather evidence iteratively. Small language models (SLMs) running locally often struggle to
control this loop. They may answer too early, retrieve after collecting enough evidence,
or ask irrelevant sub-questions.

%% [§1 ¶3]
These control decisions are small, structured prediction problems: \emph{Is the evidence
sufficient? What should happen next? How much evidence is still missing?} Using the
generator ties decision quality to model size and consumes tokens at every
step. Laya~\cite{laya} is a non-autoregressive \dq{System~1} decision model. It answers
multiple-choice, yes/no, and ordinal questions in one forward pass, returning a
probability for each option. Training uses strictly proper scoring rules, and temperature
scaling calibrates these probabilities. Most adaptive retrieval methods use the
generator's signals or decide once per query (Section~\ref{sec:related}). When external
calibrated control should replace the generator's control, and how this depends on
generator size, remain open questions.

%% [§1 ¶4]
Our $2\times2$ factorial design uses a shared decomposition loop to isolate the
controller's effect. We vary the \emph{evidence source}
(closed-book recall or RAG with dense retrieval and re-ranking) and the \emph{loop
controller} (the LLM or Laya). Laya's decisions are enforced at every hop, rather than
offered as an optional tool. Thinking mode is disabled for all models. We evaluate five
Qwen3.5 sizes~\cite{qwen35} on three multi-hop benchmarks, with 200 questions per
configuration ($5\times4\times3=60$ configurations). We report exact match (EM),
token-level F1, cost, paired bootstrap 95\% confidence intervals (CIs), and the
controller's Brier score and expected calibration error (ECE). We ask
three research questions (RQs).
\begin{itemize}
  \item \textbf{RQ1 (Accuracy).} With the decomposition loop fixed, does replacing the
        generator's control with a calibrated external decision model change multi-hop QA
        accuracy (EM and F1)? Does the effect depend on the evidence source?
  \item \textbf{RQ2 (Cost).} How does external control affect tokens, hops, end-to-end
        latency, and tool-call reliability per question?
  \item \textbf{RQ3 (Scale).} How do these effects change with generator scale, from
        0.8B to 27B parameters?
\end{itemize}

%% [§1 ¶5]
Our four contributions address these research questions as follows.
\begin{itemize}
  \item A factorial framework separating evidence and controller effects within one loop,
        with wrong answers classified at the tool, retrieval, stopping, or answering
        stage (RQ1).
  \item A weak-supervision method that simulates agent states from supporting-paragraph
        annotations to train a calibrated multi-hop controller for RQ1--RQ3.
  \item An analysis of tokens, hops, latency, and tool-call reliability under explicit
        control in local SLM agents (RQ2).
  \item An evaluation of five Qwen3.5 sizes on three datasets showing that the accuracy
        effect of external control reverses with generator scale (RQ3).
\end{itemize}

%% [§1 ¶6]
Sections~\ref{sec:related}--\ref{sec:experiments} cover related work, methods, and results.
Section~\ref{sec:discussion} answers the RQs and discusses limitations.
Section~\ref{sec:conclusion} concludes the paper.

\section{Related Work}
\label{sec:related}

%% [§2 ¶1]
\textbf{Retrieval Pipelines.} RAG~\cite{lewis2020rag} supplies external passages to a
generator. Yang \emph{et~al.}~\cite{yang2025enhancing} combined BGE-M3
embeddings~\cite{chen2024bgem3}, an exact inner-product index~\cite{douze2024faiss}, and a
BGE cross-encoder re-ranker~\cite{chen2024bgem3,bgerankerv2m3}. They reported accuracy
gains for nearly all of 17 LLMs on two Chinese QA benchmarks. We adopt this pipeline and
study who decides when to retrieve and stop.

%% [§2 ¶2]
\textbf{Multi-Hop QA and Agents.} HotpotQA~\cite{yang2018hotpotqa},
2WikiMultihopQA~\cite{ho2020twowiki}, and MuSiQue~\cite{trivedi2022musique} require facts
across Wikipedia paragraphs and label the supporting ones. The first two
use EM and token-level F1, while MuSiQue uses answer F1. We report EM and F1 on all
three and use their annotations to build corpora, derive weak labels, and classify
errors. AutoGen~\cite{wu2024autogen} and its continuation AG2~\cite{ag2} let an LLM call
tools across turns. Loops usually end through the generator's answer tool or a step
limit. Tool-call reliability varies with model size.

%% [§2 ¶3]
\textbf{Adaptive and Dynamic Retrieval.} FLARE~\cite{jiang2023flare} retrieves when an
upcoming sentence has low-confidence tokens. Adaptive-RAG~\cite{jeong2024adaptiverag}
uses query complexity to choose no, single-step, or iterative retrieval.
DRAGIN~\cite{su2024dragin} decides when and what to retrieve from the generator's
information needs. SeaKR~\cite{yao2025seakr} uses uncertainty in the generator's internal
states to trigger retrieval. These methods use the generator's signals or decide once per
query. S2G-RAG~\cite{li2026s2grag} checks evidence sufficiency at each turn.
RAG-on-a-Diet~\cite{ding2026ragdiet} trains a reinforcement-learning agent to select the
smallest model that can handle each hop and to stop at a confidence threshold.

%% [§2 ¶4]
\textbf{Calibrated Decision Models.} Given a textual state and a question,
Laya~\cite{laya} returns option probabilities in one forward pass. It supports
\texttt{choice}, \texttt{noul} (yes/no, returning $P(\text{true})$), and \texttt{score}
(ordinal) questions. Miscalibrated sufficiency probabilities can trigger early or late
stops. We therefore report the Brier score and ECE.

%% [§2 ¶5]
\textbf{Positioning.} Deciding when to retrieve and when to stop is an established
problem, and we do not claim a new decision layer for it. Our controller is calibrated,
non-autoregressive, separate from the generator, and trained only on weak labels. Our
factorial design changes only the evidence source and controller. RAG-on-a-Diet selects
among generators to reduce cost. We instead vary generator size to study how it changes
the accuracy-efficiency trade-off of external control.

\section{Method}
\label{sec:method}

\subsection{Problem Setting and $2\times2$ Factorial Design}
\label{sec:method:design}

%% [§3.1 ¶1]
For a multi-hop question $q$ from dataset $d$, the agent uses an LLM $M$ and corpus
$\mathcal{C}_d$. Its state $s=(q,\,Q_{\mathrm{sub}},\,E)$ tracks the question,
sub-questions, and collected evidence. The \emph{evidence source}
$e\in\{\cond{recall},\cond{rag}\}$ determines where evidence comes from. The \emph{loop
controller} $c\in\{\cond{llm},\cond{laya}\}$ decides when to stop and answer. Their
combinations give the four conditions in
Table~\ref{tab:conditions}.

%% [§3.1 ¶2]
\begin{table}[t]
  \centering
  \caption{The four conditions of the $2\times2$ factorial design. All conditions share
  the loop of Algorithm~\ref{alg:loop}.}
    \label{tab:conditions}
  \small
  \begin{tabular}{|l|l|l|}
    \hline
    Condition & Evidence source $e$ & Controller $c$ \\
    \hline
    Recall+LLM  & \cond{recall} (closed-book) & \cond{llm} \\
    \hline
    RAG+LLM     & \cond{rag} (retrieve + re-rank) & \cond{llm} \\
    \hline
    Recall+Laya & \cond{recall} (closed-book) & \cond{laya} \\
    \hline
    RAG+Laya    & \cond{rag} (retrieve + re-rank) & \cond{laya} \\
    \hline
  \end{tabular}
\end{table}

%% [§3.1 ¶3]
The design separates the two main effects and their interaction. Recall+LLM also uses
the loop: the LLM asks sub-questions, answers from memory, and decides when to stop.
It does not give a single-shot closed-book answer. Differences between \cond{llm} and
\cond{laya} therefore reflect who stops the loop and which tools are available, rather
than iterative decomposition.

\subsection{Shared Decomposition Loop}
\label{sec:method:loop}

%% [§3.2 ¶1]
Algorithm~\ref{alg:loop} shows the shared loop. At each hop, the controller can stop.
Otherwise, the LLM asks a sub-question and collects evidence through retrieval or recall.
The LLM answers from this evidence when the controller stops or the loop reaches
$H_{\max}=4$ hops.

%% [§3.2 ¶2]
\begin{algorithm}[t]
  \caption{Decomposition loop shared by the four conditions. Retrieved passages are
  de-duplicated by passage id and capped at $|E|\le8$.}
    \label{alg:loop}
    \raggedright% the wrapped Require line without stretched spaces

  \begin{algorithmic}[1]
    \Require question $q$; evidence source $e\in\{\cond{recall},\cond{rag}\}$;
             controller $c\in\{\cond{llm},\cond{laya}\}$; $H_{\max}=4$
    \State $Q_{\mathrm{sub}}\gets[\,]$; \ $E\gets[\,]$
    \For{$h=1,\dots,H_{\max}$}
      \State $d_h\gets\varnothing$ \Comment{no hint unless $c=\cond{laya}$}
      \If{$c=\cond{laya}$}
        \State $d_h\gets\mathrm{Laya}(\mathrm{BuildState}(q,Q_{\mathrm{sub}},E))$
        \IfThen{$d_h.p_{\mathrm{suf}}\ge0.5$}{\textbf{break}} \Comment{enforced}
      \EndIf
      \State $q_h\gets M.\mathrm{SubQuestion}(q,Q_{\mathrm{sub}},E,d_h)$
      \If{$c=\cond{llm}$ \textbf{and} $M$ called \texttt{final\_answer}}
        \State \textbf{break}
      \EndIf
      \State append $q_h$ to $Q_{\mathrm{sub}}$
      \If{$e=\cond{rag}$}
        \State $E\gets E\cup\mathrm{Rerank}(q_h,\mathrm{Dense}(q_h,20))[1{:}3]$
      \Else
        \State $E\gets E\cup\{M.\mathrm{Recall}(q_h)\}$
      \EndIf
    \EndFor
    \State \Return $M.\mathrm{Answer}(q,E)$
  \end{algorithmic}
\end{algorithm}

%% [§3.2 ¶3]
\textbf{Controller.} When $c=\cond{laya}$, Laya runs at each hop's start. If its
sufficiency probability is $p_{\mathrm{suf}}\ge0.5$, the loop ends and the LLM must
answer. Otherwise, its predicted next action and remaining hops guide the LLM, with
only \texttt{ask\_sub\_question} available. When $c=\cond{llm}$, the LLM can call
\texttt{ask\_sub\_question} or stop through \texttt{final\_answer}. We count the
stopping hop, so hops equal decisions. Each
question has a 300\,s budget.

%% [§3.2 ¶4]
\textbf{Evidence and failures.} When $e=\cond{rag}$, the dense retriever uses the
sub-question to find the top-20 passages in $\mathcal{C}_d$. The re-ranker's top-3
passages are added to $E$. When $e=\cond{recall}$, we add the LLM's closed-book answer
to the sub-question. A reply without a valid call to an available tool, including a
plain-text reply, ends the question as \texttt{invalid\_tool}. Exceeding the time budget
gives \texttt{timeout}. Other failures receive \texttt{error} or \texttt{oom}. All count
as incorrect.

\subsection{Laya as the Decision Layer}
\label{sec:method:laya}

%% [§3.3 ¶1]
At each hop, Laya answers three fixed English questions about the state: whether to stop,
what to do next, and how much evidence is missing. Each uses a different question type.

%% [§3.3 ¶2]
\begin{itemize}
  \item \texttt{sufficient} (\texttt{noul}): is the evidence sufficient for the original question? Laya returns $p_{\mathrm{suf}}=P(\text{true})$.
  \item \texttt{next\_action} (\texttt{choice}): (A) answer now; (B) form one more
        sub-question and search; (C) rephrase the query and search again.
  \item \texttt{remaining\_hops} (\texttt{score}): how many key evidence passages are
        missing? The options are 0, 1, or 2 (two or more).
\end{itemize}

%% [§3.3 ¶3]
\textbf{State construction.} \textsc{BuildState} converts the question, sub-questions,
and evidence into text within a 900-token limit, keeping only the first two sentences of
each item. It uses the same format for training and inference.

\subsection{Weak Supervision by State Simulation}
\label{sec:method:weak}

%% [§3.4 ¶1]
The datasets lack human labels for these three decisions, but their training splits
identify the supporting paragraphs. For a training question with gold paragraphs $G$ and
distractors $D$, we simulate four kinds of states: (1) \emph{question only}, $E=\emptyset$;
(2) \emph{partial gold}, a non-empty proper subset of $G$ plus 0--2 passages from $D$;
(3) \emph{all gold}, all of $G$ plus 0--2 passages from $D$, shuffled; and
(4) \emph{distractors only}.

%% [§3.4 ¶2]
Let $m=|G\setminus E|$ be the number of gold paragraphs missing from the evidence. The
labels are
\begin{align}
  \texttt{sufficient}      &= \mathbf{1}[m=0], \\
  \texttt{next\_action}    &=
    \begin{cases}
      \text{A} & \text{if } m=0,\\
      \text{C} & \text{if } E\neq\emptyset \text{ and } E\cap G=\emptyset,\\
      \text{B} & \text{otherwise},
    \end{cases} \\
  \texttt{remaining\_hops} &= \min(m,\,2).
\end{align}
Missing gold evidence calls for another decomposition step. If the evidence contains only
distractors, the query should be rephrased instead.

%% [§3.4 ¶3]
With a fixed seed, we sample 800, 100, and 100 training questions per dataset for
decision-layer training, calibration, and held-out evaluation. We target about
30{,}000 training decisions. Before the main runs, \texttt{sufficient} and
\texttt{next\_action} must each reach 0.75 accuracy and exceed the zero-shot base
checkpoint (evaluated only here) by at least 0.20. Mean absolute error (MAE) on
\texttt{remaining\_hops} must be at most 0.5.

\subsection{Evaluation Protocol}
\label{sec:method:eval}

%% [§3.5 ¶1]
Our primary measure is EM after official answer normalisation. We also report token-level
F1 (for MuSiQue, the maximum over the answer and its aliases). For each question, we record
end-to-end latency, device-wide peak GPU memory, LLM prompt and output tokens, hops, and
Laya-call latency. For the decision layer, we report accuracy, Brier score, ECE, and
latency per decision. EM differences include paired percentile-bootstrap 95\% CIs based
on 10{,}000 resamples of the 200 questions.

\subsection{Implementation Details}
\label{sec:method:impl}

%% [§3.6 ¶1]
\textbf{Models and Datasets:} We use five instruction-tuned Qwen3.5
models~\cite{qwen35} (0.8B, 2B, 4B, 9B, and 27B parameters) with thinking disabled,
an 8192-token context, greedy decoding, and seed 42. All run on one 24\,GB consumer GPU
(software, checkpoints, and licenses are in Appendix~\ref{sec:appendix}). We evaluate
HotpotQA~\cite{yang2018hotpotqa} in the distractor setting,
2WikiMultihopQA~\cite{ho2020twowiki} (2Wiki), and the answerable subset of
MuSiQue~\cite{trivedi2022musique}. For each dataset, we sample 200 test questions from the
validation split and the 800/100/100 decision-data questions from the training split.
No test questions enter decision-layer training, calibration, or held-out evaluation.
Before the main runs, each model answers 20 fixed tool-calling probe items. We retain
all models.

%% [§3.6 ¶2]
\textbf{Retrieval:} Following~\cite{yang2025enhancing}, we embed passages and
sub-questions with BGE-M3~\cite{chen2024bgem3}, retrieve passages by exact inner-product
search, and re-score them with the BGE cross-encoder re-ranker~\cite{bgerankerv2m3}.
Each dataset's corpus combines all supporting and distractor paragraphs from its 200 test
questions. We remove duplicates and use one chunk per paragraph. The embedder, re-ranker,
and Laya share the GPU with the generator.

%% [§3.6 ¶3]
\textbf{Controller:} We fine-tune one shared decision layer for all datasets from the
multilingual Laya checkpoint~\cite{laya}. Training uses Laya's reinforcement learning
with strictly proper scoring rules~\cite{laya}. We then fit one temperature per question
type on the calibration split (details are in Appendix~\ref{sec:appendix}).


% Every number in this section comes from <Data Root>/runs/main-20260930/summary.csv,
% runs/main-20260930/probe_results.json, runs/laya_eval/laya_eval_report.{md,json}
% (run_id main-20260930), paper/analysis/bootstrap_ci.csv, or paper/analysis/error_taxonomy.csv.
% tables/*.tex and figures/*.pdf are generated (figures and the older tables by
% `multihop-benchmark aggregate`, Ticket 11; the R8 panel tables by paper/analysis/merge_tables.py;
% tables/error_taxonomy.tex by paper/analysis/error_taxonomy.py) and are not edited by hand; each
% tables/<stem>.tex is a bare tabular and \tableorplaceholder / \widetableorplaceholder add the
% float, caption, and \label{tab:<stem>} (figures: \label{fig:<stem>}); see main.tex.
\section{Experimental Results}
\label{sec:experiments}

\subsection{Setup}
\label{sec:exp:setup}

%% [§4.1 ¶1]
Each of the 60 configurations ran once on 200 questions, giving 12{,}000 evaluations
(Section~\ref{sec:method:impl} and Table~\ref{tab:conditions}). EM and F1 are
percentages, and differences are percentage points (pp). One correct answer adds
0.5\,pp EM. Sections~\ref{sec:exp:main} and~\ref{sec:exp:efficiency} address RQ1
and RQ2, and paragraphs marked RQ3 examine model size. Section~\ref{sec:exp:errors}
analyses errors. Section~\ref{sec:discussion} answers the RQs.

\subsection{Function-Calling Probe and Decision-Layer Quality}
\label{sec:exp:laya}

%% [§4.2 ¶1]
Every model made valid tool calls on all 20 probe items
(Table~\ref{tab:controller_probe}, Panel~A), with mean latency from 134\,ms (0.8B)
to 836\,ms (27B). This probe used one tool and single-hop questions. It did not predict
the main results, where the 0.8B model's \texttt{invalid\_tool} rate under LLM control
was 0.290--0.550. Other models stayed at or below 0.030 in every configuration.

%% [§4.2 ¶2]
The decision layer met all thresholds in Section~\ref{sec:method:weak}
(Table~\ref{tab:controller_probe}, Panel~B). Accuracy reached 0.9271 on
\texttt{sufficient} (+0.3356 over zero-shot) and 0.8333 on \texttt{next\_action}
(+0.4859). MAE on \texttt{remaining\_hops} was 0.2013. For \texttt{sufficient}, the
Brier score fell from 0.5551 to 0.1343 and ECE from 0.2271 to 0.0626. All fitted
temperatures reached 5.0, the upper bound of the search range. A call answering all three
questions had a median latency of 12.6\,ms (zero-shot: 13.2\,ms).

%% [§4.2 ¶3]
\tableorplaceholder[\scriptsize]{controller_probe}{Function-calling probe with mean
latency per call (Panel A). Decision-layer quality on 4{,}032 held-out decisions
(Panel B), fine-tuned (FT) versus zero-shot (ZS): Brier for \texttt{sufficient} and
\texttt{next\_action}, MAE for \texttt{remaining\_hops}.}

\subsection{Main Results (RQ1)}
\label{sec:exp:main}

\begin{figure*}[t]
%% [§4.3 ¶1]
  \subfigureorplaceholder[0.3\textwidth]{em_vs_size_hotpotqa}{EM by model size, HotpotQA.}\hfill
%% [§4.3 ¶2]
  \subfigureorplaceholder[0.3\textwidth]{em_vs_size_2wiki}{EM by model size, 2Wiki.}\hfill
%% [§4.3 ¶3]
  \subfigureorplaceholder[0.3\textwidth]{em_vs_size_musique}{EM by model size, MuSiQue.}
\end{figure*}

%% [§4.3 ¶4]
Table~\ref{tab:main_results_panels} (Panel~A) reports EM and F1, and
Figs.~\ref{fig:em_vs_size_hotpotqa}--\ref{fig:em_vs_size_musique} show EM by model
size. The evidence source has the largest effect. RAG+LLM exceeds Recall+LLM for every
model and dataset, by 19.0--43.5\,pp EM on HotpotQA, 3.5--36.5\,pp on 2Wiki, and
3.5--30.0\,pp on MuSiQue. RAG+LLM with 27B performs best throughout (EM 67.5, 68.5,
and 37.0, respectively). No model exceeds 10.0 EM without retrieval on MuSiQue, the
hardest dataset.

%% [§4.3 ¶5]
\tableorplaceholder[\scriptsize]{main_results_panels}{Main results, 200 questions per cell.
Panel A: EM / F1 (\%). Panel B: controller effects on EM (pp) with paired bootstrap 95\% CIs:
$\Delta_{\mathrm{RAG}}$ = RAG+Laya minus RAG+LLM, and the interaction
$\Delta_{\mathrm{int}}$ = $\Delta_{\mathrm{RAG}}$ minus $\Delta_{\mathrm{Rec}}$, where
$\Delta_{\mathrm{Rec}}$ = Recall+Laya minus Recall+LLM (Appendix~\ref{sec:appendix:results}).}

%% [§4.3 ¶6]
\textbf{Model size (RQ3).} Under recall, Laya changes EM by $-2.0$ to $+2.5$\,pp for
the 2B--27B models, equivalent to at most five questions. The 0.8B model's EM point
estimates rise by 5.5, 8.0, and 0.5\,pp on HotpotQA, 2Wiki, and MuSiQue. Under
retrieval, the effect depends on size. RAG+Laya falls below RAG+LLM
by 3.5--15.5\,pp EM on every dataset for the 4B, 9B, and 27B models. The largest
drops occur on 2Wiki ($-15.5$, $-14.5$, and $-12.5$\,pp, respectively). For the 2B
and 0.8B models, RAG+Laya gives the highest EM point estimates on every dataset,
exceeding RAG+LLM by 1.5--9.0\,pp. On HotpotQA, EM rises from 40.0 to 47.5 for 2B
and from 23.0 to 30.5 for 0.8B. The interaction point estimate is negative for all nine
combinations of larger models and datasets, and positive for all six with smaller
models. F1 changes in the same direction.

%% [§4.3 ¶7]
% CI numbers from paper/analysis/bootstrap_ci.csv (paper/analysis/bootstrap_ci.py).
\textbf{Sampling uncertainty.} Under RAG, the CI for Laya's effect excludes zero in
eleven of the fifteen combinations of models and datasets
(Table~\ref{tab:main_results_panels}, Panel~B). These include 0.8B on all three
datasets (lower bounds of 2.0--4.5\,pp), seven of the nine combinations with larger
models, and 2B only on HotpotQA ($[3.0, 12.5]$). Interaction CIs exclude zero in six
of nine combinations with larger models, all negative, and two of six with smaller
models, both positive (2B on HotpotQA and 0.8B on MuSiQue).
Appendix~\ref{sec:appendix:results} lists intervals including zero and the controller
effect under recall and the retrieval effect.

\subsection{Error Analysis}
\label{sec:exp:errors}

%% [§4.4 ¶1]
% Counts and shares from paper/analysis/error_taxonomy.csv (dataset=all rows pool the three
% datasets; share = count / wrong answers of the model and condition).
Table~\ref{tab:error_taxonomy} classifies wrong answers under retrieval by stage using
gold supporting paragraphs (rules in Appendix~\ref{sec:appendix:results}). Categories
cover failed calls, complete evidence, premature stops, and retrieval misses after four
hops. Under RAG+Laya, premature stops are the most frequent
category for every model: 194--230 questions, or 45.4--65.0\% of wrong answers, compared
with 0--56 under RAG+LLM. Under RAG+LLM, the 4B--27B models most often fail despite
having retrieved all gold evidence (47.9--58.5\%). The 2B model mainly fails to retrieve
the evidence within $H_{\max}$ (315 questions, 73.1\%). The 0.8B model mainly fails
through retrieval misses (247) and invalid tool calls (211). Laya reduces retrieval
misses from 315 to 78 for 2B and from 247 to 108 for 0.8B. However, premature stopping
becomes the main source of error for every model.

%% [§4.4 ¶2]
Appendix~\ref{sec:appendix:results} gives the counts per dataset and a representative
premature stop.

%% [§4.4 ¶3]
\tableorplaceholder[\scriptsize]{error_taxonomy}{Wrong answers under retrieval by stage,
pooled over the three datasets (600 questions per row): invalid tool call (Inv.), timeout
(T/O), error (Err.), retrieval miss after all hops (Miss), all gold evidence retrieved
(Cov.), premature stop with gold evidence missing (Stop), and other (Oth.).}

\subsection{Efficiency (RQ2)}
\label{sec:exp:efficiency}

%% [§4.5 ¶1]
\widetableorplaceholder[\scriptsize]{efficiency_panels}{Efficiency per model and condition,
one panel per dataset: end-to-end p50 latency (ms), device-wide peak GPU memory (VRAM, MB),
mean prompt (In) and output (Out) tokens per question, mean hops, and mean latency per Laya
call (ms).}

\begin{figure*}[t]
%% [§4.5 ¶2]
  \subfigureorplaceholder[0.31\textwidth]{latency_vs_size}{End-to-end p50 latency (log scale)
  by model size.}\hfill
%% [§4.5 ¶3]
  \subfigureorplaceholder[0.45\textwidth]{hops_distribution}{Hops per question by condition,
  pooled over models and datasets; the stopping hop is counted.}
\end{figure*}

%% [§4.5 ¶4]
RAG+Laya uses fewer prompt tokens, output tokens, and hops than RAG+LLM throughout
(Table~\ref{tab:efficiency_panels}). Mean prompt tokens fall from 2{,}776--6{,}067
to 1{,}565--2{,}329, and mean hops from 2.67--4.00 to 2.36--2.94. Under its own
control with retrieval, 2B almost always reaches $H_{\max}$ (mean hops of 3.98, 4.00,
and 4.00). For 0.8B--9B, RAG+Laya lowers p50 latency on every dataset
(Fig.~\ref{fig:latency_vs_size}), for example from 1880 to 657\,ms for 2B on HotpotQA.
Laya calls for these models average 16.0--20.8\,ms across conditions. All one-hop
outcomes under LLM control are \texttt{invalid\_tool} failures, not deliberate stops
(Fig.~\ref{fig:hops_distribution}). RAG+Laya most often stops at hop~2.
Appendix~\ref{sec:appendix:results} reports recall costs and p95 latency.

%% [§4.5 ¶5]
\textbf{Model size (RQ3).} The 27B model is the exception. Its Laya calls took
361.3--1261.4\,ms on average. RAG+Laya had higher p50 latency than RAG+LLM on HotpotQA
(13482 versus 11918\,ms) and 2Wiki (11490 versus 11044\,ms). With the 27B model,
retrieval components, and Laya sharing the 24\,GB GPU, peak memory use reached
23{,}526--24{,}461\,MB. These slow calls are consistent with memory pressure rather
than the decision model's usual latency (a held-out median of 12.6\,ms). Memory
pressure therefore complicates latency comparisons involving Laya at 27B.

% Numbers only from summary.csv, probe_results.json, laya_eval_report.md,
% paper/analysis/bootstrap_ci.csv (CIs), and paper/analysis/error_taxonomy.csv of
% run main-20260930; run-process facts (resumes, timeouts, the single error)
% come from manifest.json and the run log and carry no result numbers.
% Each RQ is answered in the same order: answer, evidence, explanation, boundary, implication.
\section{Discussion}
\label{sec:discussion}

\subsection{Answers to the Research Questions}
\label{sec:disc:rq}

%% [§5 ¶1]
\textbf{RQ1 (Accuracy).} Laya changed EM mainly under retrieval, with effects depending
on model size. Under recall, EM changed by $-2.0$ to $+2.5$\,pp for 2B--27B,
suggesting that available evidence limits performance more than loop control.
The 0.8B point estimates rose by up to 8.0\,pp. Under retrieval, RAG+Laya exceeded
RAG+LLM by 1.5--9.0\,pp in EM point estimates for 2B and 0.8B, but fell below it
by 3.5--15.5\,pp for 4B--27B
(Table~\ref{tab:main_results_panels}). CIs excluded zero for 0.8B on every dataset,
seven of nine combinations with larger models, and 2B only on HotpotQA. Stopping helps
explain this difference. For every model, RAG+Laya most often failed by stopping with
gold evidence missing. Under LLM control, most errors followed complete evidence or the hop limit
(Section~\ref{sec:exp:errors}). The 0.8B gain also coincided with a different tool
protocol: its \texttt{invalid\_tool} rate was 0.290--0.550 with LLM control and two
tools, but 0.000 with Laya and one tool. This does not explain the 2B gain, as its rate
was at most 0.015. Held-out controller quality (\texttt{sufficient} accuracy 0.9271,
ECE 0.0626) was measured on simulated states. The 2B gain remains uncertain on 2Wiki
and MuSiQue. Loop control should be
evaluated together with the evidence source and tool protocol.

%% [§5 ¶2]
\textbf{RQ2 (Cost).} With retrieval, external control reduced tokens and hops in every
configuration and lowered p50 latency for 0.8B--9B. Mean prompt tokens fell from
2{,}776--6{,}067 to 1{,}565--2{,}329 (Table~\ref{tab:efficiency_panels}). Laya calls
averaged 16.0--20.8\,ms for these models across conditions.
Apart from the 0.8B model's invalid calls, tool calls were reliable, with one
\texttt{error} and no \texttt{oom}. Earlier stopping explains the savings. RAG+Laya
most often stopped at hop~2 (Fig.~\ref{fig:hops_distribution}), without an extra LLM
call for the stopping decision. The LLM still generated the answer. Under its own
control with retrieval, 2B almost always reached $H_{\max}$. At 27B,
memory pressure complicates latency comparisons (run accounting is in
Appendix~\ref{sec:appendix}). These consistent savings justify external control only
where accuracy is maintained (RQ3).

%% [§5 ¶3]
\textbf{RQ3 (Scale).} External control had opposite accuracy effects across model sizes,
but saved tokens and hops throughout. Interaction point estimates were negative in
all nine combinations with the three larger models and positive in all six with the
two smaller models. CIs excluded zero in six of nine negative and two of six positive
cases. Under its own control with retrieval, 2B almost never chose to stop. The
4B--27B models used more hops than under Laya (for example,
3.19--3.42 versus 2.53--2.55 on 2Wiki) and achieved higher EM. These runs do not
identify a mechanism. For 4B--27B, Laya stopped with gold evidence missing on 194--230
questions that were then answered incorrectly. Possible
causes include the fixed threshold of 0.5 with temperatures at the search limit (5.0),
differences between simulated and retrieved states, and bridging facts lost through
\textsc{BuildState}'s truncation. Smaller models' gains are consistent with one available
tool and difficulty judging when to stop. External control may help small models manage
the loop while ending stronger models' searches too early. This reversal was observed
with one model family, loop, and threshold. Benefits should be evaluated at each size
rather than assumed to carry over.

\subsection{Limitations}
\label{sec:disc:limits}

%% [§5 ¶4]
With thinking disabled, we cannot compare Recall+LLM with a thinking-enabled
baseline. Each corpus contains paragraphs from only 200 test questions, rather than full
Wikipedia, potentially making retrieval easier than in open-domain settings. Each
configuration ran once with one seed. Bootstrap intervals capture variation across
questions, not runs. Differences up to 2\,pp (four questions) should not be
over-interpreted. We used only Qwen3.5 models and English datasets, excluding translated
datasets, cloud models, human evaluation, a zero-shot Laya condition, and Laya as a
passage re-ranker. Supporting-paragraph annotations supplied weak labels to train one
shared Laya on 800 questions per dataset. All results depend on the shared loop. Gold
paragraphs serve as a proxy for sufficient evidence, although answers may not require
all of them. All fitted temperatures reached the search limit (5.0), so the reported
ECE does not establish optimal calibration.

% Numbers only from summary.csv and laya_eval_report.md of run main-20260930.
\section{Conclusion}
\label{sec:conclusion}

%% [§6 ¶1]
We studied when calibrated external control should replace the generator's control of
multi-hop retrieval. The controller trained on weak supporting-paragraph labels met
all pre-registered thresholds, with \texttt{sufficient} accuracy of 0.9271 and ECE of
0.0626. Retrieval contributed most, improving EM by up to 43.5\,pp. \textbf{RQ1.}
External control changed EM mainly under retrieval. Errors most often involved stopping
before all gold evidence was retrieved. \textbf{RQ2.} With retrieval, Laya reduced prompt
tokens, output tokens, and hops throughout. It eliminated invalid tool calls for 0.8B
and lowered p50 latency for 0.8B--9B. \textbf{RQ3.} Under retrieval, Laya raised EM
point estimates by 1.5--9.0\,pp for 2B and 0.8B. The 4B--27B models achieved
3.5--15.5\,pp higher EM under their own control. CIs support the 0.8B gain throughout
and the 2B gain only on HotpotQA.

%% [§6 ¶2]
External control did not benefit every model. For the smallest models, it coincided with
no invalid tool calls (0.8B), shorter retrieval loops, and higher EM point estimates at
lower cost. For larger models, it stopped earlier, often before bridging evidence
arrived, saving tokens and hops but lowering EM. This accuracy-efficiency trade-off
depended on generator size and should be evaluated at each scale.


% ieeenat_fullname.bst (CVPR kit): full author names; reference titles are not quoted. Normal size
% (10 pt), as in the reference paper (R9 round 2).
\bibliographystyle{ieeenat_fullname}
\bibliography{refs}

% Appendix (R9 round 2): after the References, from a new page (page 9), outside the 8-page budget.
% cvpr.sty numbers the appendix sections "A.", "B.". Main-text tables are cited in the main text;
% the appendix tables are cited only in the appendix (paper/check_paper.sh checks both).
\clearpage
\appendix
% Appendix (revision R9 round 2): after the References, from page 9, outside the 8-page budget.
% A: implementation and reproducibility details (facts from manifest.json, spike_results.json, and
% laya_eval_report.md; seed from paper/analysis/bootstrap_ci.csv). B: supplementary results; every
% number comes from paper/analysis/error_taxonomy.csv, paper/analysis/bootstrap_ci.csv, summary.csv,
% or the records of run main-20260930. The appendix tables are generated (error_taxonomy.py,
% merge_tables.py) and cited only here; the main text cites the appendix sections, not these tables.
% Identifiers are \texttt with \allowbreak after '/' (R9 round 3: \path spread them out in loose lines).
\section{Implementation and Reproducibility Details}
\label{sec:appendix}

%% [§A ¶1]
\textbf{Software and hardware.} Ollama~0.34.4 serves the language models (tags
\texttt{qwen3.5:0.8b}, \texttt{:2b}, \texttt{:4b}, \texttt{:9b}, \texttt{:27b}) on an
NVIDIA GeForce RTX 4090 under Ubuntu on WSL2 with Python~3.12. The embedder is
\texttt{BAAI/\allowbreak bge-m3}, and the re-ranker is
\texttt{BAAI/\allowbreak bge-reranker-\allowbreak v2-m3}. A FAISS flat inner-product
index runs on the CPU. Because the 27B model ran fully on the GPU, BGE-M3, the re-ranker,
and Laya also remained on the GPU in all 60 configurations. We obtain the datasets from
Hugging Face (MuSiQue as \texttt{musique-ans}) and convert them to a common schema.
AG2~1.1.0 could not disable thinking in our setting, so we implement the loop directly
with the Ollama Python client.

%% [§A ¶2]
\textbf{Decision-layer training.} Each simulated state presents a passage both as a full
paragraph and as its supporting sentences (MuSiQue keeps the full paragraph). This exposes
the layer to retrieved and recall-like evidence. We fine-tune the decision layer from
\texttt{convaiinnovations/\allowbreak laya-multilingual} (322M parameters, input length 1024 tokens) on
32{,}451 simulated decisions and calibrate it on 3{,}996, in 10.8\,min.
Laya's reinforcement learning~\cite{laya} adds Gaussian noise to option logits for
exploration. A strictly proper scoring rule rewards each sample against its weak label.
Advantages use a group baseline, and a soft cross-entropy term guides training. We then
fit one temperature per question type on the calibration split. All three reached
5.0, the upper end of the search range (Section~\ref{sec:disc:limits}).

%% [§A ¶3]
\textbf{Records and post-processing.} For each question, we log the prediction,
correctness, per-step trace (sub-questions, retrieved passage ids, Laya outputs), latency,
tokens, and peak GPU memory. All results come from these records, the probe output, and
the Laya evaluation report. Post-processing scripts reproduce the paired bootstrap
intervals (seed 20261001), error taxonomy, and all result tables with identical bytes
on every run. Questions that run out of GPU memory are retried once with a 4096-token
context and recorded as \texttt{oom} if they fail again. Licenses are CC~BY-SA~4.0 for
HotpotQA, CC~BY~4.0 for MuSiQue, MIT for BGE-M3, and Apache-2.0 for 2WikiMultihopQA,
Laya, Qwen3.5, and the BGE re-ranker.

%% [§A ¶4]
\textbf{Run accounting.} The single \texttt{error} (9B, 2Wiki, RAG+LLM) was a malformed
tool-call XML reply. A question answered after $k$ sub-questions has $k+1$ hops unless
it reaches $H_{\max}$. Laya never stopped at the first hop, when evidence was empty.
Failed questions count hops only up to failure, lowering the 0.8B model's mean under
LLM control. Only two questions reached the 300\,s budget (2Wiki, 27B, RAG+LLM).
They timed out three times in total, one of them twice. Both eventually finished after
extra retries that no other question received. Thus, the final \texttt{timeout\_rate}
is zero in every configuration.


\section{Supplementary Results}
\label{sec:appendix:results}

%% [§B ¶1]
% Rules of paper/analysis/error_taxonomy.py (docstring), checked in this order.
\textbf{Error taxonomy rules.} Each wrong answer under retrieval receives one category,
checked in the following order. A failed call has final status \texttt{invalid\_tool},
\texttt{timeout}, or \texttt{error}. Evidence covered means that every gold supporting
paragraph was retrieved at some hop. A premature stop leaves a gold paragraph missing
when the last step stops the loop (Laya's $p_{\mathrm{suf}}\ge0.5$ or the LLM's
\texttt{final\_answer}). A retrieval miss leaves a gold paragraph missing even though
the last step still retrieved. Remaining cases are classified as other.
Table~\ref{tab:error_taxonomy_by_dataset} reports counts per dataset, and
Table~\ref{tab:error_taxonomy} pools them.

%% [§B ¶2]
\tableorplaceholder[\scriptsize]{error_taxonomy_by_dataset}{Wrong answers under retrieval by
stage, per dataset (200 questions per row). Columns follow
Table~\ref{tab:error_taxonomy}.}

%% [§B ¶3]
% Case from records hotpotqa__qwen3.5-27b__rag_{laya,llm}.jsonl, question_id
% 5abb2f2f554299232ef4a3d7 (gold paragraphs: Harry Emerson Fosdick, Riverside Church).
\textbf{A representative premature stop} (HotpotQA, 27B). The question was \dq{Harry
Emerson Fosdick was called to serve as pastor at a church that opened its doors on what
date?} The first hop retrieved the paragraph on Fosdick but not the one on Riverside
Church. Laya then judged the evidence sufficient ($p_{\mathrm{suf}}=0.983$), and the
model answered that the information was unavailable. Under its own control, the same
model next asked \dq{When did Riverside Church open its doors?} and answered correctly
(October 5, 1930).

%% [§B ¶4]
% CI numbers from paper/analysis/bootstrap_ci.csv.
\textbf{Intervals including zero; recall and retrieval effects.} Under RAG, the CI for the
Laya effect (Table~\ref{tab:main_results_panels}, Panel~B) includes zero for 2B on 2Wiki
($[-3.5, 6.5]$) and MuSiQue ($[0.0, 8.0]$), 9B on HotpotQA, and 27B on MuSiQue.
Table~\ref{tab:effects_recall_retrieval} lists the
controller effect under recall ($\Delta_{\mathrm{Rec}}$ = Recall+Laya minus Recall+LLM) and
the retrieval effect ($\Delta_{\mathrm{Ret}}$ = RAG+LLM minus Recall+LLM) with paired
bootstrap 95\% CIs. The $\Delta_{\mathrm{Rec}}$ CI excludes zero only for 0.8B on HotpotQA
and 2Wiki. The $\Delta_{\mathrm{Ret}}$ CI excludes zero in every pair except 0.8B on 2Wiki
($[-1.0, 8.0]$).

%% [§B ¶5]
\tableorplaceholder[\scriptsize]{effects_recall_retrieval}{Recall effect
$\Delta_{\mathrm{Rec}}$ and retrieval effect $\Delta_{\mathrm{Ret}}$ on EM (pp), estimate
[paired bootstrap 95\% CI], per model and dataset.}

%% [§B ¶6]
\textbf{Costs under recall and tail latency.} Recall+Laya also uses fewer prompt tokens
than Recall+LLM (Table~\ref{tab:efficiency_panels}). Its stopping decision requires no
extra LLM call, which partly explains this reduction. Under recall, p50 latency falls on
every dataset only for 0.8B and 2B (for example, 1277 versus 2733\,ms for 2B on
HotpotQA). Table~\ref{tab:efficiency_p95} reports end-to-end p95 latency alongside the
p50 results in Table~\ref{tab:efficiency_panels}.

%% [§B ¶7]
\tableorplaceholder[\scriptsize]{efficiency_p95}{End-to-end p95 latency (ms) per model,
condition, and dataset.}



\end{document}
